\centerline{\Huge Условия и решения задач.}
\centerline{\bf \Huge 5 класс}
 
\problem{5.1}На доске записано выражение: *** + *** = 950. Замените звёздочки цифрами так, чтобы равенство стало верным, и все девять цифр на доске были различными (цифры 9, 5 и 0 нельзя использовать вместо звёздочек).

\hspace{3mm}

Решение:\\
Возможные ответы:\\
214 + 736 = 950\\
216 + 734 = 950\\
234 + 716 = 950\\
236 + 714 = 950\\
Дополнительных объяснений не требуется.

\hspace{3mm}

Критерии:\\
Приведён любой из возможных ответов: \textbf{7 баллов.}\\
Приведёно верное равенство, в котором какие-то две цифры совпадают: \textbf{2 балла.} \\
Получение неверного равенства или равенства с совпадением более двух цифр или многократным совпадением более двух и более цифр: \textbf{0 баллов.}\\


\problem{5.2}В каждой из двадцати школ города Элисты обучаются поровну учеников. Если в каждой школе будет обучаться на 50 учеников меньше, то во всех школах останется столько учеников, сколько их раньше было суммарно в девятнадцати школах. Сколько учеников изначально обучалось в каждой школе?

\hspace{3mm}

Ответ: 1000 учеников.

Решение:\\
\textit{I вариант}: \\
Всего во всех школах стало обучаться на 20 · 50 = 1000 учеников меньше. До этого было 20 школ с одинаковым количеством учеников, а стало 19 школ с одинаковым количеством учеников. Значит, общее количество учеников уменьшилось на одну школу и в то же время на 1000 учеников. Следовательно, в каждой школе обучается по 1000 учеников.

\newpage
\hspace{10mm}

\textit{II вариант}: Пусть в каждой школе обучается по X учеников. Тогда составим уравнение:\\$20\cdot X - 20\cdot50 = 19\cdot X\\
X = 20\cdot50 = 1000$. \\Следовательно, в каждой школе обучается по 1000 учеников.

\hspace{3mm}

Критерии:\\
Любое верное решение не подбором и с верным обоснованием: \textbf{7 баллов.}\\
В целом верное решение с верным обоснованием, но с допущением арифметических ошибок: \textbf{5 баллов.}\\
Решение подбором с обоснованием подбора: \textbf{4 балла.}\\
Идеи, продвигающие ход решения, но недосаточные для получения верного ответа: \textbf{2 балла.}\\
Неверное решение или идеи, не продвигающие ход решения: \textbf{0 баллов.} \\
Ответ без обоснования: \textbf{0 баллов.}

\hspace{3mm}

\problem{5.3}Три друга - Алдар, Бадма и Хонгор закончили решать городской этап олимпиады по математике, и каждый воскликнул. Алдар: «Я точно умнее Бадмы!», Бадма: «Алдар определенно не самый умный из нас!», Хонгор: «Эх, Алдар умнее меня…». Известно, что только самый умный солгал, а остальные сказали правду. Может ли самым умным быть Алдар?

\hspace{3mm}

Ответ: Нет, не может.

Решение:\\
Предположим, что Алдар самый умный. Значит, по условию задачи он солгал, а остальные сказали правду. Но Бадма сказал, что Алдар не самый умный из нас, а он самый умный по предположению. Следовательно, Бадма лжёт, однако лгать может только Алдар. Противоречие. Значит, наше предположение о том, что Алдар самый умный, было неверным. Отсюда следует, что Алдар не самый умный.  

\hspace{3mm}

Критерии:\\
Полностью верное решение со всеми обоснованиями: \textbf{7 баллов.}\\

\newpage
\hspace{3mm}

\noindentИдеи, продвигающие ход решения, но недосаточные для получения верного ответа: \textbf{2 балла.}\\
Неверное решение  или идеи, не продвигающие ход решения: \textbf{0 баллов.}\\
Ответ без обоснования: \textbf{0 баллов.}

\begin{wrapfigure}{R}{0.3\textwidth}
\vspace{-20mm}
\centering
\includegraphics[width=0.35\textwidth]{Городская Олимпиада Элиста 2023/Условия/picture_5.5.jpg}
\caption{Рис. 1}
\end{wrapfigure}

\hspace{3mm}

\problem{5.4} Однажды Пюрвя вырезал по линиям клетки из клетчатого листа бумаги 4 одинаковые фигурки и решил их склеить друг с другом. У него получилась замечательная фигура (рисунок справа). Определите, как разрезать замечательную фигуру на изначальные 4 одинаковые фигурки Пюрви?

\hspace{3mm}

\begin{wrapfigure}{R}{0.3\textwidth}
\vspace{0mm}
\centering
\includegraphics[width=0.35\textwidth]{Городская Олимпиада Элиста 2023/Решения/picture_5.5_sol.jpg}
\caption{Рис. 1}
\end{wrapfigure}

\hspace{3mm}

Решение:\\
Одно из разрезаний: рисунок справа снизу\\
Дополнительных объяснений не требуется.

Критерии:\\
Любое верное разрезание, возможно, отличающееся от предложенного в решении: \textbf{7 баллов.}\\
Получена площадь фигурки, равная 5 клеткам: \textbf{3 балла.}\\
Получена площадь замечательной фигуры, равная 20 клеткам: \textbf{1 балл.}\\
Любое неверное разрезание (не по линиям, не из четырех фигур, не из одинаковых фигур и т.д.): \textbf{0 баллов.}

\hspace{3mm}

\problem{5.5}На шахматной доске необычного размера 2023 клетки в длину и 2023 клетки в ширину, состоящей из 2046264 белых полей и 2046265 чёрных, (расположенных в шахматном порядке) стоит по одной белой пешке на каждом белом поле, а на чёрных полях ничего не стоит. В какой-то момент времени Рома решил одновременно взять каждую белую пешку и переставить в какое-то соседнее по стороне поле. 

\newpage
\hspace{3mm}

\noindentКакое минимальное количество чёрных полей будет не занято, если на одном поле может стоять более одной пешки?

\hspace{3mm}

Ответ: 1 черное поле.

Решение:\\
Заметим, что соседние с белой клеткой по стороне только чёрные клетки. Следовательно, Рома переставляет каждую белую пешку с белого поля только на чёрное. Всего белых клеток 2046264, значит, и белых пешек на них стоит 2046264, а чёрных полей 2046265. Следовательно, по обратному принципу Дирихле: клеток, то есть чёрных полей, больше, чем кроликов, то есть белых пешек. Таким образом, найдется хотя бы одна клетка, то есть хотя бы одно чёрное поле, которое будет пустым. А значит, минимальное количество чёрных полей, которые будут не заняты, равно единице. (Можно доказывать без обратного принципа Дирихле через метод от противного)


\hspace{3mm}

Критерии:\\
Полностью верное решение со всеми обоснованиями: \textbf{7 баллов.}\\
Идеи, продвигающие ход решения, но недосаточные для получения верного ответа: \textbf{2 балла.}\\
Неверное решение или идеи, не продвигающие ход решения: \textbf{0 баллов.}\\
Ответ без обоснования: \textbf{0 баллов.}

